\section{Introduction}\label{sec:intro}
A model can recognize patterns in machine-written text without using the same patterns to produce its own writing. Distinguishing these two possibilities matters for interpreting representation-based detectors: an accurate linear readout does not by itself identify a controllable writing mechanism. Even when an intervention lowers an AI-style score, it may do so by damaging coherence, removing facts, or exploiting weaknesses in the evaluator.

\para{Can reading out AI style predict control of new writing?}
Linear probes and mean-difference detectors extract substantial human--machine information from frozen language-model activations \cite{quaremba2026,vishnyakov2026}. Prior sparse-autoencoder work also steers individual stylistic features, with factual side effects at strong doses \cite{kuznetsov2025}. The unresolved question is therefore more specific than whether style can be steered: can a human--AI readout direction change independent judgments of new text while preserving coherence and content, and can this effect be distinguished from nuisance style and persona factors?

We test this conjunction in paired pretrained and instruction-tuned \qwen checkpoints. We fit directions to human and GPT-4o continuations from matched source documents, select a layer on validation readout accuracy, and intervene during generation on 64 unseen factual records. \Figref{fig:method} summarizes the protocol. Signed doses, projection ablation, equal-norm random directions, nuisance controls, and a plain human-style prompt separate several explanations for a score change. A frozen RoBERTa detector and a blinded Gemini judge evaluate the outputs; explicit noninferiority margins test preservation.

Negative steering lowers judged AI-likeness by 20.04/100 points, but reduces coherence by 0.591/5 points and fails both preservation gates. The encoder change is $+0.061$, with a confidence interval spanning zero. Strong readout in both backbones also fails to generalize to Llama-3 base continuations. We make three contributions:
\begin{itemize}[leftmargin=*,itemsep=0pt,topsep=0pt]
\item We test whether a source-matched human--AI residual readout controls new generations using signed, random, ablation, nuisance, and prompt interventions.
\item We quantify detector disagreement and quality costs with paired inference, calibration, and prespecified preservation margins.
\item We compare domain and generator readout transfer across paired base and instruction backbones, documenting a base-generator transfer failure and no observed amplification in AUROC.
\end{itemize}
\Secref{sec:related} positions the study; \secref{sec:method} defines the protocol; \secref{sec:results} reports the results; and \secref{sec:discussion} examines their limits.
\begin{figure}[t]
\centering
\setlength{\fboxsep}{6pt}
\fbox{\begin{minipage}{0.27\linewidth}\small\textbf{1. Fit and select}\\Matched human/GPT-4o texts\\Train: 384 sources\\Validation: 96 sources\\Last-token layer 24 direction\end{minipage}}
\hfill$\rightarrow$\hfill
\fbox{\begin{minipage}{0.27\linewidth}\small\textbf{2. Intervene}\\64 new factual records\\Signed additions / ablation\\Equal-norm random controls\\Nuisance and prompt controls\end{minipage}}
\hfill$\rightarrow$\hfill
\fbox{\begin{minipage}{0.27\linewidth}\small\textbf{3. Evaluate jointly}\\Blinded AI-style judge\\Frozen encoder detector\\Coherence and fact coverage\\Unsupported claims\end{minipage}}
\caption{Readout-to-generation test. Layer selection uses held-out readout accuracy, without generation scores. We pair generation outcomes by factual record and test style alongside preservation. A style-score change alone does not satisfy the hypothesis.}\label{fig:method}
\end{figure}

